\chapter{Scénář skupinových diskusí}
\label{app:fgd_scenario}

Níže je uveden scénář moderátora pro skupinové diskuse. Celková délka diskuse: 80 minut.

\section*{Blok I — Vnímání produktu a~první dojem (20 min)}

\begin{enumerate}
  \item „Když se řekne mražený dezert, co si vybavíte jako první?"
  \item „Znáte značku Franui? V~čem je podle vás jiná než běžné nanuky nebo mražené dřeně?"
  \item „Jaké emoce ve vás vyvolává představa mraženého ovoce v~čokoládě?"
  \item „Koupili byste si takový produkt? Proč ano / proč ne?"
\end{enumerate}

\section*{Blok II — Senzorické testování (30 min)}

\textit{Respondentům jsou rozdány vzorky produktu Berrie.}

\begin{enumerate}
  \item „Jaký máte pocit z~textury po 5, 10, 15 minutách od vytažení z~mrazáku? Je ovoce stále pevné a~osvěžující, nebo už působí příliš měkkým dojmem?"
  \item „Jak vnímáte poměr mezi ovocem a~čokoládou? Není čokoláda příliš dominantní, nebo naopak málo výrazná oproti kyselosti ovoce?"
  \item „Jak na vás působí kombinace mraženého ovoce a~čokoládové krusty?"
  \item „Kterou variantu (borůvky, jahody, banán) byste si vybrali při opakovaném nákupu?"
\end{enumerate}

\section*{Blok III — Design a~positioning (30 min)}

\begin{enumerate}
  \item „Působí na vás obal produktu jako prémiová pochutina, kterou byste si koupili na Rohlík.cz nebo v~supermarketu?"
  \item „Na základě ochutnávky a~vzhledu balení, jakou cenu byste u~tohoto produktu očekávali? Působí na vás 150g balení jako adekvátní porce?"
  \item „Věřili byste českému původu ovoce, i~když je produkt v~mraženém stavu? Jak moc je pro vás informace o~lokálnosti a~BIO kvalitě důležitá?"
  \item „Zaujala by vás v~budoucnu více veganská varianta, nebo rozšíření o~jiné druhy ovoce?"
\end{enumerate}

% ============================================================================
\chapter{Struktura dotazníku CAWI}
\label{app:questionnaire}

\textit{Kompletní dotazník bude vložen po finalizaci výzkumného nástroje.}

\section*{Sekce 1 — Demografické údaje}
\begin{itemize}
  \item Věk, pohlaví, region, vzdělání
  \item Využívání sociálních sítí (Instagram, TikTok)
\end{itemize}

\section*{Sekce 2 — Chuťové preference}
\begin{itemize}
  \item „Které z~následujících druhů ovoce byste obalené v~čokoládě vyzkoušeli nejraději? (maliny, borůvky, banán, jahody, mix)"
  \item Preference typu čokolády (bílá, mléčná, hořká)
\end{itemize}

\section*{Sekce 3 — Nákupní chování}
\begin{itemize}
  \item „Při jaké příležitosti byste si mražené ovoce v~čokoládě nejpravděpodobněji koupili?"
  \item Preferovaný distribuční kanál
  \item Frekvence potenciálního nákupu
\end{itemize}

\section*{Sekce 4 — Cenová citlivost (PSM)}
\begin{itemize}
  \item „Při jaké ceně byste produkt považovali za tak drahý, že byste o~jeho koupi vůbec neuvažovali?"
  \item „Při jaké ceně byste produkt považovali za tak levný, že byste pochybovali o~jeho kvalitě?"
  \item „Jaká cena je podle vás už vysoká, ale stále byste o~koupi uvažovali?"
  \item „Jaká cena je podle vás výhodná (skvělý nákup)?"
\end{itemize}

% ============================================================================
\chapter{Rozpočet výzkumu}
\label{app:budget}

\begin{table}[H]
\centering
\caption{Detailní rozpočet marketingového výzkumu značky Berrie}
\label{tab:budget_detail}
\begin{tabularx}{\textwidth}{Xlr}
\toprule
\textbf{Položka} & \textbf{Specifikace} & \textbf{Částka (Kč)} \\
\midrule
\multicolumn{3}{l}{\textbf{Online dotazníkové šetření (CAWI)}} \\
\quad Přístup k panelu respondentů & 400 respondentů & 35\,000 \\
\quad Nastavení a programování dotazníku & & 7\,000 \\
\midrule
\multicolumn{3}{l}{\textbf{Skupinové diskuse (FGD)}} \\
\quad Pronájem diskusní místnosti & 2 sessions × 80 min & 8\,000 \\
\quad Práce moderátora & & 8\,000 \\
\quad Odměny pro účastníky & 16 osob × 375 Kč & 6\,000 \\
\midrule
\multicolumn{3}{l}{\textbf{Senzorické testování (CLT)}} \\
\quad Pronájem prostor & & 6\,000 \\
\quad Práce zaškolených tazatelů & & 10\,000 \\
\midrule
\multicolumn{3}{l}{\textbf{Suroviny a vzorky}} \\
\quad Výroba vzorků Berrie & & 6\,000 \\
\quad Nákup konkurenčních produktů (Franui) & & 4\,000 \\
\midrule
\textbf{Celkem} & & \textbf{90\,000} \\
\bottomrule
\end{tabularx}
\end{table}
